\scanpage{72}
% archive foliation: 71
de $X_i$ suivant les points à l'infini\ldots).

Cl. que $\widehat X_i$ est automatiquement $\omega$-orienté, donc aussi
\[
\widehat X=\coprod\widehat X_i
\]
De cette façon, la catégorie envisagée (des surfaces topologiques ``spéciales'' $\omega$-orientées \inserted{avec homéomorphismes $\omega$-orientés}) devient équivalente à celle des surfaces topologiques \inserted{$\omega$-orientées} à comp. conn. compactes \uline{munies} d'une partie \deleted{\unsure{f}} $S$ discrète, avec les isom. de couples $(\widehat X,S)$.

De $\overline{\mathrm m}$, la catégorie des couples \deleted{$\unsure{(X,S)|(X,o,S)}$} $(X,o,\unsure{\mathcal J|J})$ avec $(X,o)$ comme dessus ($o$ la $\omega$-orientation) et $\mathcal J$ partie discrète de $X$, est équiva-\\
lente à celle des systèmes $(\widehat X,o,S,\mathcal J)$, où $(\widehat X,o)$ est une surface $\omega$-orientée à comp. conn. compactes,
